\cotefile{157-5}
\scanpage{453}\mspage{365}
% archive sheet number: 452
est \unsure{aussi|encore} bien dans $\mathrm{Fl}\,K'$.

Ainsi, $K'$ est une sous-catégorie de $J$, mais est-elle filtrante ? La condition $\mathrm{PS}\,1'$ est évidente, puisque $e$ est un plus petit élément de $J$, reste la condition $\mathrm{PS}\,2$ de la double flèche. On voit aisément que cette condition est satisfaite dans $K'$, sss elle est satisfaite pour les double-flèches
\diagram{\begin{tikzcd}[ampersand replacement=\&,column sep=large]
e \arrow[r,shift left=.48ex,"p^{\unsure{r|s}}"] \arrow[r,shift right=.48ex,"p^s"'] \& e \&[-1em] \text{avec}\quad r,s\geq 0\unsure{,|.}
\end{tikzcd}}

(\deleted{\unsure{ou encore}}En effet, \deleted{\unsure{des flèches|deux flèches}} On se ramène aussitôt au cas de \unsure{deux|ttes} flèches de but $e$, et on note que \unsure{tte} flèche de $K'$ de but $e$ s'écrit comme un composé
\diagram{\begin{tikzcd}[ampersand replacement=\&,column sep=small]
k\arrow[r,"u"] \& e\arrow[r,"p^r"] \& e\arrow[r,"v"] \& \unsure{e}
\end{tikzcd}}, avec $r\geq0$ et $u,v$
des flèches dans $K'$ \unsure{\ldots|.} En composant
avec \diagram{\begin{tikzcd}[ampersand replacement=\&,column sep=small]e\arrow[r,"p"] \& e\end{tikzcd}}, on trouve alors, comme $pv=p$
\unsure{:} \diagram{\begin{tikzcd}[ampersand replacement=\&,column sep=small]
k\arrow[r,"u"] \& e\arrow[r,"p^{r+1}"] \& e,\quad\text{soit}\quad \& k\arrow[r,"{\unsure{fu|pu}}"] \& e\arrow[r,"p^r"] \& e
\end{tikzcd}})

Mais on voit aussitôt que la possibilité d'égaliser $(p^r,p^s)$ dans $K'$ équivaut à l'existence de $t\unsure{>|\geq}0$ tel que $p^{r+t}=p^{s+t}$. Donc $K'$ est filtrante sss \deleted{\unsure{$\varphi$}} $r\longmapsto p^r$ est une \deleted{\unsure{suite}} stationnaire (faire $r\unsure{=|\geq}0$, $s=1$, d'où $p^t=p^{t+1}$ pour \unsure{grand|$t$ grand}). Ce n'est
